\documentclass{article}
\usepackage{tcee-repaso}
\tcetema[Gobernanza económica de la UE]{3.B.45}{La gobernanza económica de la Unión Europea y de la zona euro. El Semestre Europeo. Las reglas fiscales y el Procedimiento de Desequilibrios Macroeconómicos. El Mecanismo Europeo de Estabilidad. La respuesta común de política fiscal a partir de la covid-19.}
\tcefecha{19/11/2026}
\begin{document}
\cabecerarepaso
\begin{columnas}

% ==========================================================================
\bloque[2,5]{Introducción}
\begin{itemize}
\item \textbf{Enganche}: «quien bien te quiere te hará llorar».
  \begin{itemize}
  \item Las reglas obligan a ajustar: los excesos de uno los pagan todos; la ayuda común llega con condiciones.
  \item MEDE a España (2012): \dato{41.333 M€}{MEDE} con MoU; NGEU (2020) a cambio de reformas.
  \item ¿Cómo se gobierna una moneda común sin un Tesoro común?
  \end{itemize}
\item \textbf{Relevancia}:
  \begin{itemize}
  \item 21 Estados en el euro (Bulgaria, 2026); zona del euro: \dato{déficit 3,0\,\%, deuda 87,7\,\%}{Eurostat, 2025}.
  \item España: senda de gasto hasta 2031; \dato{deuda 100,5\,\%}{2025} \vertema{4.B.26}.
  \item Reforma de 2024; cláusula nacional de defensa (2025).
  \end{itemize}
\item \textbf{Contextualización}: \concepto{gobernanza económica}: reglas, procedimientos e instituciones para coordinar y vigilar las políticas fiscales y estructurales.
  \begin{itemize}
  \item UEM asimétrica: monetaria única (BCE) \vertema{3.B.43}; fiscal nacional coordinada (art.~121).
  \item 3 piezas: prevención (reglas, PDM, Semestre); gestión de crisis (MEDE); respuesta común (SURE, NGEU, SAFE).
  \end{itemize}
\item \textbf{Problemática}: ¿por qué reglas y por qué fallaron? ¿Qué cambia en 2024? ¿Quién ayuda y a cambio de qué? ¿Hay ya política fiscal común?
\item \textbf{Estructura}: \hyperref[rep:I]{I} fundamentos; \hyperref[rep:II]{II} reglas, PDM y Semestre; \hyperref[rep:III]{III} crisis y respuesta común.
\end{itemize}
\esquemacante{%
\apartadocante{Introducción}{2,5}
\apartadocante{I. Fundamentos e instituciones}{7,5}
\subcante{1. ¿Por qué reglas? (sesgo deficitario)}{3}
\subcante{2. Marco jurídico e institucional}{2}
\subcante{3. De Maastricht a 2024}{2,5}
\apartadocante{II. Reglas fiscales, PDM y Semestre}{11}
\subcante{1. Reglas fiscales}{7,5}
\subcante{2. El PDM}{2}
\subcante{3. El Semestre Europeo}{1,5}
\apartadocante{III. Gestión de crisis y respuesta común}{8}
\subcante{1. El MEDE}{3}
\subcante{2. La respuesta desde la covid-19}{2,5}
\subcante{3. ¿Capacidad fiscal central?}{2,5}
\apartadocante{Conclusión}{1}}

% ==========================================================================
\bloque[7,5]{I. Fundamentos e instituciones\label{rep:I}}
\begin{ideaclave}
En una unión monetaria el coste del déficit de cada Estado lo pagan en parte los demás: cada uno se endeuda de más. De ahí las reglas, el no rescate (art.~125) y la vigilancia común, que se han endurecido y flexibilizado a golpe de crisis.
\end{ideaclave}

\sub{I.1. ¿Por qué necesita reglas una unión monetaria?\label{rep:I.1}}
\subsub{a) Los argumentos}
\begin{itemize}
\item \emph{Externalidades fiscales} \vertema{3.B.16}, por 3 vías:
  \begin{itemize}
  \item \emph{tipo de interés}: quien se endeuda sube el de todos; el BCE puede endurecer;
  \item \emph{inflación}: tentación de monetizar si se espera que el BCE compre deuda \vertema{3.A.39};
  \item \emph{contagio financiero}: a bancos y deuda de los demás (2010-2012).
  \end{itemize}
\item \emph{Sesgo deficitario}: ilusión fiscal, deuda estratégica, fondo común \vertema{3.A.39}; en la unión, mayor.
\item \emph{Riesgo moral}: \concepto{no corresponsabilidad} (\emph{no bail-out}, art.~125); sin financiación del BCE (art.~123).
\item \emph{La disciplina de mercado no basta}: Grecia al tipo de Alemania hasta 2008; pánico en 2010-2012.
  \begin{itemize}
  \item Deuda en una moneda que no emiten: crisis que se cumplen solas (\autor{De Grauwe}, 2011).
  \end{itemize}
\end{itemize}
\subsub{b) El sesgo deficitario en una unión monetaria}
\begin{itemize}
\item \emph{Idea}: problema de polizón sobre un bien común, la estabilidad de la moneda (\autor{Beetsma}-\autor{Uhlig}, 1999; \autor{Chari}-\autor{Kehoe}, 2007).
\end{itemize}
\subsubsub{i) Supuestos}
\begin{itemize}
\item $n$ Estados con un banco central; beneficio $v(b_i)=ab_i-b_i^2/2$.
\item El banco central acomoda: $\pi=\theta\bar b$; coste $\frac\lambda2\pi^2$; equilibrio de \autor{Nash}.
\end{itemize}
\subsubsub{ii) Desarrollo}
\begin{itemize}
\item \emph{Nash}: cada uno ve $\theta/n$: $a-b=\lambda\theta^2b/n$.
\item \emph{Unión} (planificador): $a-b=\lambda\theta^2b$:
\funcionclave{b^{N}=\frac{a}{1+\lambda\theta^2/n}>b^{C}=\frac{a}{1+\lambda\theta^2}}
  \begin{itemize}
  \item $a=6$, $\lambda\theta^2=2$, $n=4$: $b^N=4$, $b^C=2$.
  \item Con $n=1$ coinciden: el sesgo adicional nace de compartir moneda.
  \end{itemize}
\end{itemize}
\graficorepaso[0.75\linewidth]{sesgo-union-monetaria}{Cada Gobierno ve solo $1/n$ del coste de su deuda}
\subsubsub{iii) Implicaciones}
\begin{itemize}
\item Reglas comunes en torno a $b^C$: es la justificación del 3\,\% y del PEC.
\item La independencia del BCE no basta (\autor{Kydland}-\autor{Prescott}, 1977): por eso arts.~123-125.
\item Reglas creíbles: el incentivo a incumplir crece con $n$.
\end{itemize}
\subsubsub{iv) Evidencia y valoración}
\begin{itemize}
\item 2003 (Alemania y Francia); 24 PDE (2009-2011); previsiones optimistas (\autor{Frankel}-\autor{Schreger}, 2013); Grecia 2009.
\item Explica por qué reglas, no cuáles ni cómo hacerlas cumplir.
\end{itemize}

\sub{I.2. Marco jurídico e institucional\label{rep:I.2}}
\subsub{a) Base en los Tratados (art.~5 TFUE: coordinación)}
\begin{itemize}
\item 119: principios; 121: coordinación y \emph{supervisión multilateral} (brazo preventivo y PDM).
\item 126 + Protocolo 12: \concepto{3\,\%} de déficit y \concepto{60\,\%} de deuda.
\item 123-125: sin financiación monetaria, acceso privilegiado ni rescate.
\item 136: zona del euro; 136.3 (2011): MEDE. 122: MEEF, SURE, NGEU.
\end{itemize}
\subsub{b) Instituciones}
\begin{itemize}
\item \concepto{Comisión}: vigila y propone. \concepto{Ecofin}: decide por mayoría cualificada.
  \begin{itemize}
  \item Art.~126: sin el afectado; 126.9 y 126.11, solo la zona del euro (139.4).
  \item Sanciones: \concepto{mayoría cualificada inversa} (2011).
  \end{itemize}
\item \concepto{Eurogrupo} (Protocolo 14): \autor{Pierrakakis} desde 12-12-2025; \concepto{Cumbre del Euro} (\autor{Costa}).
\item \concepto{Parlamento}: colegisla; diálogo económico.
\item IFI: Consejo Fiscal Europeo (2016, \autor{Hasekamp}); AIReF (2013, \autor{Olóndriz}) \vertema{4.B.21}.
\end{itemize}
\subsub{c) Grados de coordinación}
\begin{itemize}
\item Monetaria \emph{única}; fiscal \emph{estrecha} (reglas, sanciones); estructural \emph{débil} (recomendaciones, método abierto).
\end{itemize}
\graficorepaso[0.9\linewidth]{arquitectura-gobernanza}{Orienta, propone, decide; prevenir, corregir, gestionar crisis}

\sub{I.3. De Maastricht a 2024\label{rep:I.3}}
\begin{enumerate}
\item \emph{Maastricht} (1992): criterios de convergencia; arts.~121 y 126.
\item \emph{PEC} (1997): 1466/97 (\concepto{preventivo}, OMP) y 1467/97 (\concepto{correctivo}, PDE, multas hasta 0,5\,\%).
\item \emph{Reforma de 2005}, tras el choque de 2003: OMP por país, «otros factores», plazos más largos.
\item \emph{Crisis del euro}:
  \begin{itemize}
  \item \concepto{Six-Pack} (2011): Semestre, regla de gasto, 1/20, mayoría cualificada inversa, PDM;
  \item \concepto{TECG} (2012): saldo estructural $\geq-0{,}5$\,\%; art.~135 CE (2011), LO 2/2012;
  \item \concepto{Two-Pack} (2013): presupuestos antes del 15-10.
  \end{itemize}
\item \emph{2015-2024}: Cinco presidentes y flexibilidad; cláusula general (2020-2023); marco de 2024.
\end{enumerate}

% ==========================================================================
\bloque[11]{II. Reglas fiscales, PDM y Semestre\label{rep:II}}
\begin{ideaclave}
El PEC fue complejo, procíclico e incumplido. Desde 2024: 3\,\% y 60\,\%, pero el instrumento es una senda de gasto neto pactada con cada Estado a partir de un ASD; el PDM vigila el resto y el Semestre lo ordena.
\end{ideaclave}

\sub{II.1. Las reglas fiscales\label{rep:II.1}}
\subsub{a) El PEC hasta 2024 y sus problemas}
\begin{itemize}
\item \emph{Preventivo}: OMP estructural ($\geq-1$\,\%; $-0{,}5$ TECG), 0,5 puntos/año; regla de gasto \vertema{4.B.4}.
\item \emph{Correctivo}: PDE por déficit ($>$3\,\%) o deuda (1/20 del exceso).
\item \emph{4 problemas}:
  \begin{enumerate}
  \item complejo (vademécum de $>$100 páginas);
  \item no observable (brecha de producción);
  \item procíclico (2011-2013);
  \item incumplido: ninguna multa (2016: cero a ES y PT); 1/20 imposible.
  \end{enumerate}
\end{itemize}
\subsub{b) ¿Qué regla? Déficit frente a gasto}
\begin{itemize}
\item \emph{Idea}: reglas frente a discreción \vertema{3.A.36}; de saldo, deuda o gasto: ¿cuál deja actuar a los estabilizadores?
\item \emph{Supuestos}: $x=u+m\hat g$; ingresos $\tau x$. Déficit: $\hat g=\tau x$; gasto: $\hat g=0$.
\item \emph{Desarrollo}:
\funcionclave{x^{D}=\frac{u}{1-m\tau}>x^{E}=u}
  \begin{itemize}
  \item $m=1$, $\tau=0{,}45$: la de déficit amplifica por 1,8.
  \end{itemize}
\item \emph{Implicaciones}:
  \begin{itemize}
  \item saldo total, procíclica; estructural, no observable;
  \item gasto: anticíclica y observable, pero necesita un ancla (ASD).
  \end{itemize}
\item \emph{Extensiones}: descontar medidas de ingresos; cláusulas de escape; estándares (\autor{Blanchard}-\autor{Leandro}-\autor{Zettelmeyer}, 2021).
\item \emph{Evidencia}: las de gasto se cumplen más y son menos procíclicas (\autor{Cordes} \emph{et al.}, 2015).
\end{itemize}
\graficorepaso[0.7\linewidth]{regla-deficit-gasto}{La regla de gasto deja actuar a los estabilizadores}
\subsub{c) El marco de 2024}
\begin{itemize}
\item Reglamentos 2024/1263 (preventivo) y 2024/1264 (correctivo); Directiva 2024/1265 (marcos nacionales); en vigor el 30-4-2024.
\item \emph{No cambian}: 3\,\% y 60\,\%; PDE por déficit. \emph{Cambian}, 6 piezas:
  \begin{enumerate}
  \item \emph{ASD} y \concepto{trayectoria de referencia}: deuda a la baja 10 años tras el ajuste.
  \item \concepto{Gasto primario neto}: sin intereses, desempleo cíclico, fondos UE ni medidas de ingresos.
  \item \emph{Planes} de 4-5 años; ajuste en 4, hasta 7 con reformas e inversiones.
  \item \emph{Salvaguardias}: deuda $-1$ punto/año si $>$90\,\%, $-0{,}5$ si 60-90\,\%; déficit estructural hasta 1,5\,\% (0,4 puntos/año; 0,25 en 7 años).
  \item \emph{Cuenta de control}: desviación $>$0,3 anual o 0,6 acumulada $\Rightarrow$ \concepto{PDE por deuda}.
  \item \emph{Correctivo}: 0,5 puntos/año; multas de 0,05\,\% del PIB por semestre.
  \end{enumerate}
\item \emph{Cláusulas}: \concepto{general} (recesión de la UE) y \concepto{nacional} (circunstancias excepcionales).
\end{itemize}
\graficorepaso[0.85\linewidth]{ciclo-marco-2024}{ASD $\to$ plan $\to$ senda $\to$ informe $\to$ cuenta de control $\to$ PDE}
\subsub{d) La sostenibilidad de la deuda (ASD)}
\begin{itemize}
\item \emph{Supuestos}: $d_t=\frac{1+i}{1+\gamma^N}d_{t-1}-s_t$ \vertema{3.A.39}; ingresos crecen con $\gamma^N_p$; $s=\tau-e$.
\item \emph{Desarrollo}: la Comisión fija el $\Delta s$ anual (escenarios adversos y estocásticos); $\Delta e\approx e(\gamma_{GN}-\gamma^N_p)$:
\funcionclave{\gamma_{GN}\approx\gamma^N_p-\frac{\Delta s}{e}}
  \begin{itemize}
  \item 4,5\,\% $-$ 0,4/0,41 $\approx$ 3,5\,\%: la senda de España en 2026.
  \end{itemize}
\item \emph{Simulación} (deuda 100,5\,\%):
  \begin{itemize}
  \item sin ajuste, baja hasta 2029 ($r<g$) y sube: 108\,\% en 2041;
  \item plan a 4 años (0,6 puntos): 73\,\%; a 7 años (0,35): 76\,\%.
  \end{itemize}
\end{itemize}
\graficorepaso[0.7\linewidth]{simulacion-deuda}{El plan a 7 años suaviza el ajuste}
\begin{itemize}
\item \emph{Implicaciones}: una regla por país; más tiempo a cambio de reformas; el envejecimiento cuenta.
\item \emph{Evidencia}: plan de España: 96,6\,\% (2028), 90,6\,\% (2031), 76,8\,\% (2041).
\item \emph{Valoración}: más económico y sencillo, pero supuestos discutibles, negociación bilateral y poca capacidad de hacerse cumplir.
\end{itemize}
\subsub{e) La aplicación en 2025-2026}
\begin{itemize}
\item \emph{Planes}: 21 sendas el 21-1-2025 (no los Países Bajos); en 7 años: BE, DE, ES, FR, IT, AT, RO, FI.
\item \emph{PDE}: 7 en julio de 2024 + Rumanía (2020); Austria (2025), Finlandia y Bulgaria (2026); Malta cerrado (2026).
  \begin{itemize}
  \item \concepto{10 abiertos} (nov. 2026): AT, BE, BG, FI, FR, HU, IT, PL, RO, SK; ninguno por deuda.
  \end{itemize}
\item \emph{Cláusula de defensa} (\emph{ReArm Europe}): hasta 1,5 puntos/año en 2025-2028.
  \begin{itemize}
  \item 15 Estados (julio de 2025) + Alemania, Austria y España (6-2026); seguridad energética, 0,3.
  \end{itemize}
\end{itemize}
\graficorepaso[0.7\linewidth]{pde-abiertos}{24 en 2010; 0 en 2019; 10 en 2026}
\subsubsub{España \vertema{4.B.23}}
\begin{itemize}
\item PDE de 2009-2019; 2025: déficit 2,4\,\%, deuda 100,5\,\%.
\item \emph{Plan 2025-2028}, ajuste en 7 años: gasto neto 3,7/3,5/3,2/3,0\,\%; déficit 2,5 $\to$ 1,6\,\%.
\item \emph{Cumplimiento}:
  \begin{itemize}
  \item 2025: 4,8\,\% (desviación 0,4 puntos, dentro de la cláusula);
  \item AIReF (nov. 2026): 7,3\,\% en 2026 y 5,1\,\% en 2027: incumplimiento.
  \end{itemize}
\end{itemize}
\pareja{deficit-deuda}{España: deuda 119,3\,\% (2020) $\to$ 100,5\,\% (2025)}{gasto-neto-espana}{2025: 4,8\,\% frente a 3,7\,\%; 2026: 7,3\,\%}

\sub{II.2. El PDM\label{rep:II.2}}
\begin{itemize}
\item \emph{Por qué}: la crisis nació en desequilibrios privados y exteriores (ES e IE con superávit en 2007) \vertema{3.B.26}; Six-Pack (Reglamentos 1176/2011 y 1174/2011).
\item \emph{Fases}:
  \begin{enumerate}
  \item \emph{alerta}: informe con un cuadro de indicadores;
  \item \emph{examen exhaustivo}: sin desequilibrios, desequilibrios o excesivos;
  \item \emph{corrección}: procedimiento de desequilibrio excesivo, plan de acción; multa de hasta 0,1\,\% (ZE).
  \end{enumerate}
\item \emph{Cuadro} (alertas, no reglas; asimétrico): cuenta corriente $-4$/$+6$\,\%; PIIN $-35$\,\%; CLU $+9$\,\% (ZE); vivienda $+9$\,\%; deuda privada, pública 60\,\%; paro 10\,\%.
\item \emph{2026}: 7 exámenes; EL, NL, SE sin; IT, HU, SK con; RO excesivos.
\item \emph{España}: desequilibrios desde 2012 (excesivos en 2013); sin ellos desde 2024 \vertema{4.A.29}.
\item \emph{Valoración}: mejor análisis, poco cumplimiento; el correctivo nunca se ha abierto; superávits alemanes sin medidas.
\end{itemize}

\sub{II.3. El Semestre Europeo\label{rep:II.3}}
\begin{itemize}
\item Desde 2011: calendario común para que las recomendaciones lleguen \emph{antes} de los presupuestos.
\item \emph{Ciclo}: noviembre, \emph{paquete de otoño}; 30-4, informes de progreso; junio, \emph{paquete de primavera} y \concepto{REP}; julio, el Consejo las adopta; 15-10, presupuestos (ZE).
\item \emph{MRR} (2021-2026): las REP con dinero detrás \vertema{3.B.37}; sin él, se cumplen poco.
\end{itemize}
\graficorepaso[\linewidth]{semestre-europeo}{Recomendaciones antes de los presupuestos}

% ==========================================================================
\bloque[8]{III. Gestión de crisis y respuesta común\label{rep:III}}
\begin{ideaclave}
El MEDE presta con condiciones; ante las crisis comunes, la UE se endeuda en común (SURE, NGEU, SAFE). Falta una capacidad fiscal permanente.
\end{ideaclave}

\sub{III.1. El MEDE\label{rep:III.1}}
\subsub{a) Del préstamo a Grecia al MEDE}
\begin{itemize}
\item 2010, sin instrumento y con el art.~125; con el FMI \vertema{3.B.22}:
  \begin{itemize}
  \item préstamos bilaterales a Grecia (52.900 M€);
  \item \concepto{MEEF} (UE, art.~122.2; 60.000, presupuesto);
  \item \concepto{FEEF} (sociedad de los Estados del euro; 440.000).
  \end{itemize}
\item Art.~136.3 (2011); Tratado del MEDE (2-2-2012), operativo el 8-10-2012.
\end{itemize}
\subsub{b) Qué es y cómo funciona}
\begin{itemize}
\item \emph{Intergubernamental} (Luxemburgo); 21 miembros (Bulgaria, 29-6-2026).
\item Junta de Gobernadores (ministros), Consejo de Administración, director \autor{Gramegna}; mutuo acuerdo; urgencia 85\,\%.
\item Capital 714.700 M€ (81.700 desembolsado); capacidad \concepto{500.000} (432.000 disponibles).
\item \emph{Instrumentos}: préstamos con programa; líneas precautorias; compras de deuda; recapitalización indirecta y directa (nunca usada).
\item \emph{Condicionalidad}: MoU (troika); puerta de las OMT \vertema{3.B.43}. Acreedor preferente tras el FMI.
\end{itemize}
\subsub{c) Los programas}
\begin{itemize}
\item Grecia (3; $\approx$256.600 M€; quita privada 2012; fin 2018); Irlanda y Portugal (FEEF y MEEF); Chipre (6.300, \emph{bail-in}) \vertema{3.B.44}; España.
\item Línea pandémica (2\,\% del PIB, 240.000): nadie la pidió (estigma).
\end{itemize}
\graficorepaso[0.75\linewidth]{rescates-paises}{Grecia, 256.600 M€; España, 41.300 (vivo: 7.300)}
\subsubsub{El préstamo a España (2012) \vertema{4.A.28}}
\begin{itemize}
\item Cajas (Bankia); prima $>$600 pb. Hasta 100.000 M€ (9-6-2012); MoU 20-7-2012.
\item \concepto{41.333} al FROB (recapitalizar y Sareb); condiciones solo financieras; salida en enero de 2014.
\item Vivos $\approx$7.300 (último vencimiento, 2027).
\end{itemize}
\subsub{d) Reforma y futuro}
\begin{itemize}
\item Reforma (2021): \concepto{respaldo común} del FUR (68.000) \vertema{3.B.44}; bloqueada por Italia (2023).
\end{itemize}

\sub{III.2. La respuesta desde la covid-19\label{rep:III.2}}
\begin{itemize}
\item Perturbación común, sin riesgo moral: lo contrario de 2010 \vertema{3.A.40}.
\item 4 redes (2020): MEDE, SURE, BEI (25.000), NGEU; y el PEPP del BCE.
\end{itemize}
\subsub{a) Cláusula general de salvaguardia}
\begin{itemize}
\item 23-3-2020 a 2023; déficit de la zona del euro: 0,5\,\% (2019) $\to$ 7,0\,\% (2020).
\end{itemize}
\subsub{b) El SURE (Reglamento 2020/672, art.~122)}
\begin{itemize}
\item Préstamos con \emph{bonos sociales} para ERTE: \concepto{98.400 M€ a 19}; España, 21.324.
\item 31,5 M de trabajadores; 1,5 M de parados evitados; 9.000 M€ de intereses ahorrados.
\item No es transferencia; embrión de reaseguro de desempleo \vertema{3.B.40}.
\end{itemize}
\graficorepaso[0.6\linewidth]{sure-paises}{SURE: España, 21.324 M€ (1,9\,\% del PIB)}
\subsub{c) El NGEU \vertema{3.B.37}}
\begin{itemize}
\item 750.000 M€ (2018) de deuda de la UE, a devolver hasta 2058; MRR.
\item \emph{Salto}: deuda para \emph{transferir}; redistribución entre Estados; ligado al Semestre.
\item 450.000 pagados (79\,\%) al 1-10-2026; PIB hasta +1,4\,\% (2026).
\item \emph{Límite}: temporal y excepcional (Constitucional alemán, 2022).
\end{itemize}
\subsub{d) REPowerEU y SAFE}
\begin{itemize}
\item \emph{REPowerEU} (2022): préstamos del MRR sin pedir (225.000).
\item \emph{ReArm Europe} (800.000): cláusula + \concepto{SAFE} (2025/1106): 150.000 en préstamos; 19 Estados (Polonia 43.700; España 1.000).
\end{itemize}

\sub{III.3. ¿Capacidad fiscal central?\label{rep:III.3}}
\begin{itemize}
\item \emph{Idea}: sin tipo de cambio ni tipo de interés propios, transferencias (\autor{Kenen}) \vertema{3.B.16}; externalidad de demanda agregada (\autor{Farhi}-\autor{Werning}, 2017).
\item \emph{Modelo}: $u_A=-u_B$; $T_i=-\phi x_i$; $x_i=u_i+mT_i$:
\funcionclave{x_i=\frac{u_i}{1+m\phi}}
  \begin{itemize}
  \item $T_A=-T_B$: se financia sola; amortigua $m\phi/(1+m\phi)$ (un tercio con $m=1$, $\phi=0{,}5$).
  \end{itemize}
\item \emph{Implicaciones}: pequeña estabiliza mucho; ante golpes comunes, deuda común. \emph{Extensiones}: riesgo moral; unión de transferencias.
\item \emph{Evidencia}: EE.~UU. (\autor{Asdrubali}-\autor{Sørensen}-\autor{Yosha}, 1996): mercados 39\,\%, crédito 23\,\%, federal 13\,\%.
\item \emph{Valoración}: \autor{Draghi} (750.000-800.000 M€/año), \autor{Letta} (2024); MFP 2028-2034: Fondo de Competitividad (409.000) y mecanismo de crisis (395.000), sin estabilización automática.
\end{itemize}
\graficorepaso[0.7\linewidth]{capacidad-fiscal}{Las transferencias amortiguan un tercio del golpe}

% ==========================================================================
\bloque[1]{Conclusión}
\begin{itemize}
\item \textbf{Recapitulación}: reglas porque cada Estado se endeuda de más; marco de 2024 (gasto neto y ASD); PDM y Semestre; MEDE con condiciones; deuda común ante crisis comunes.
\item \textbf{Relevancia}: para España, ancla y restricción (2012, NGEU, senda hasta 2031) \vertema{4.B.26}.
\item \textbf{Extensiones}: \vertema{3.B.16} \vertema{3.B.43} \vertema{3.B.44} \vertema{3.B.37} \vertema{3.A.39} \vertema{4.B.23}.
\item \textbf{Opinión}: mejor marco, credibilidad en juego (10 PDE, incumplimientos); sin capacidad fiscal, la próxima crisis dependerá de una cumbre.
\item \textbf{Idea final}: «quien bien te quiere te hará llorar»: reglas que obligan y ayuda con condiciones.
  \begin{itemize}
  \item Sin cumplimiento, llanto sin disciplina; ante un golpe común, la ayuda puede llegar antes y sin castigo.
  \end{itemize}
\end{itemize}

% ==========================================================================
\anexorepaso
\bloque{Datos clave}
\begin{itemize}
\item Zona del euro: \dato{déficit 3,0\,\%, deuda 87,7\,\%}{2025}; España: \dato{2,4\,\% y 100,5\,\%}{2025}.
\item \dato{10 PDE abiertos}{nov.\ 2026}; nunca una multa.
\item MEDE: \dato{500.000 M€}{capacidad}; España, 41.333.
\item SURE: \dato{98.400 M€ a 19}{CE}; MRR: \dato{450.000 pagados (79\,\%)}{1-10-2026}.
\end{itemize}
\bloque{Autores y fechas}
\hito{1977}{\autor{Kydland}-\autor{Prescott}}
\hito{1992}{Maastricht; PEC (1997)}
\hito{1996}{\autor{Asdrubali}-\autor{Sørensen}-\autor{Yosha}}
\hito{1999}{\autor{Beetsma}-\autor{Uhlig}: sesgo}
\hito{2011}{Six-Pack; TECG (2012)}
\hito{2012}{MEDE; préstamo a España}
\hito{2017}{\autor{Farhi}-\autor{Werning}}
\hito{2020}{Cláusula general; SURE; NGEU}
\hito{2024}{Nuevo marco; \autor{Draghi}}
\bloque{Preguntas del tribunal}
\begin{enumerate}
\item \textbf{¿Por qué no basta la disciplina de mercado?} Diferenciales nulos hasta 2008 y pánico después.
\item \textbf{¿Qué es el gasto primario neto?} Sin intereses, desempleo cíclico, fondos UE ni medidas de ingresos.
\item \textbf{¿Por qué una regla de gasto?} Observable y anticíclica.
\item \textbf{¿Es el MEDE una institución de la UE?} No: intergubernamental (136.3).
\item \textbf{¿Por qué el NGEU no viola el art.~125?} La UE se endeuda para un fin común (art.~122).
\end{enumerate}
\begin{trampa}
\begin{itemize}
\item Salvaguardia: \emph{1} punto si $>$\emph{90\,\%}; \emph{0,5} entre 60 y 90.
\item Cuenta de control: \emph{0,3}/\emph{0,6}; multas \emph{0,05\,\%} por semestre.
\item MEEF: \emph{UE} (art.~122); FEEF: sociedad; MEDE: intergubernamental.
\end{itemize}
\end{trampa}

\end{columnas}
\end{document}
